\documentclass[11pt]{article}
\usepackage{amsmath,amssymb,amsthm}
\usepackage{algorithm}
\usepackage{algorithmic}
\usepackage{fullpage}
\usepackage{enumitem}
\usepackage{hyperref}
\newtheorem{theorem}{Theorem}
\newtheorem{lemma}{Lemma}
\newtheorem{assumption}{Assumption}
\newcommand{\E}{\mathbb{E}}
\newcommand{\R}{\mathbb{R}}

\begin{document}

\section*{Randomized Coordinate Descent with Probability‑Weighted Updates (RCD‑PWU)}

\section*{Mathematical Formulation}
Let $f:\R^{d}\rightarrow\R$ be a (possibly non‑convex) differentiable function.  
For each coordinate $i\in\{1,\dots ,d\}$ let $p_i>0$ be a probability with
\[
\sum_{i=1}^{d}p_i=1 .
\]
Denote by $e_i\in\R^{d}$ the $i$‑th canonical basis vector and by
\[
\nabla_i f(w)=\frac{\partial f(w)}{\partial w_i}
\]
the $i$‑th partial derivative.  

At iteration $t$ a coordinate $i_t$ is drawn independently from the
distribution $\{p_i\}_{i=1}^{d}$.  Define the *scaled stochastic gradient*
\[
g_t \;:=\; \frac{1}{p_{i_t}}\; \nabla_{i_t} f(w_t)\, e_{i_t}\in\R^{d}.
\tag{1}
\]
Note that $\E[g_t\mid w_t]=\nabla f(w_t)$.  
With a stepsize $\eta>0$ the update reads
\[
w_{t+1}=w_t-\eta\, g_t .
\tag{2}
\]
Equivalently, only the sampled coordinate is modified:
\[
w_{t+1}^{(i)}=
\begin{cases}
w_t^{(i)}-\displaystyle\frac{\eta}{p_i}\,\nabla_i f(w_t), &\text{if }i=i_t,\\[1ex]
w_t^{(i)}, &\text{otherwise.}
\end{cases}
\tag{3}
\]

\section*{Pseudocode}
\begin{algorithm}[H]
\caption{RCD‑PWU}
\begin{algorithmic}[1]
\STATE \textbf{Input:} stepsize $\eta>0$, probability vector $p\in\Delta_{d}$,
initial point $w_0\in\R^{d}$, maximum iterations $T$.
\FOR{$t=0,\dots ,T-1$}
    \STATE Sample $i_t\sim\{p_i\}_{i=1}^{d}$.
    \STATE Compute partial gradient $g_{t}^{(i_t)}=\nabla_{i_t}f(w_t)$.
    \STATE Form scaled gradient $g_t=\frac{1}{p_{i_t}} g_{t}^{(i_t)}e_{i_t}$.
    \STATE Update $w_{t+1}=w_t-\eta g_t$.
\ENDFOR
\STATE \textbf{Output:} $\{w_t\}_{t=0}^{T}$ (or an averaged iterate).
\end{algorithmic}
\end{algorithm}

\section*{Dry Run Example}
Consider the one‑dimensional function $f(w)=(w-3)^2$.
Here $d=1$, $p_1=1$, $\nabla f(w)=2(w-3)$ and we choose $\eta=0.1$,
$w_0=0$.

\begin{center}
\begin{tabular}{c|c|c|c}
$t$ & $w_t$ & $\nabla f(w_t)$ & $f(w_t)$\\ \hline
0 & $0$ & $2(0-3)=-6$ & $9$\\
1 & $w_1 = 0-0.1\cdot\frac{1}{1}(-6)=0.6$ &
       $2(0.6-3)=-4.8$ & $(0.6-3)^2=5.76$\\
2 & $w_2 = 0.6-0.1(-4.8)=1.08$ &
       $2(1.08-3)=-3.84$ & $(1.08-3)^2=3.6864$\\
3 & $w_3 = 1.08-0.1(-3.84)=1.464$ &
       $2(1.464-3)=-3.072$ & $(1.464-3)^2=2.3635$\\
\end{tabular}
\end{center}
The objective value strictly decreases, illustrating the algorithm.

\newpage

\section*{Assumptions}
\begin{assumption}[Coordinate‑wise $L_i$‑smoothness]\label{ass:smooth}
For each $i\in\{1,\dots ,d\}$ there exists $L_i>0$ such that for all
$w\in\R^{d}$ and any scalar $\alpha$,
\[
\bigl|\nabla_i f(w+\alpha e_i)-\nabla_i f(w)\bigr|\le L_i|\alpha|.
\tag{4}
\]
\end{assumption}

\begin{assumption}[Lower boundedness]\label{ass:bound}
The function $f$ is bounded below: $f^{\inf}:=\inf_{w}f(w)>-\infty$.
\end{assumption}

\begin{assumption}[Finite variance of stochastic gradient]\label{ass:var}
Define the random vector $g_t$ in \eqref{1}.  There exists $\sigma^2<\infty$ such that
\[
\E\!\left[\|g_t-\nabla f(w_t)\|^2\mid w_t\right]\le \sigma^2 .
\tag{5}
\]
\end{assumption}

\begin{assumption}[Stepsize]\label{ass:stepsize}
The constant stepsize satisfies
\[
0<\eta\le \frac{1}{\max_{i}L_i}.
\tag{6}
\]
\end{assumption}

\section*{Main Convergence Theorem}
\begin{theorem}[Expected gradient norm decay]\label{thm:main}
Under Assumptions \ref{ass:smooth}–\ref{ass:stepsize}, after $T\ge1$
iterations of RCD‑PWU we have
\[
\frac{1}{T}\sum_{t=0}^{T-1}\E\!\bigl[\|\nabla f(w_t)\|^2\bigr]
\;\le\;
\frac{2\bigl(f(w_0)-f^{\inf}\bigr)}{\eta\,T}
\;+\;
2\eta\,\sigma^2 .
\tag{7}
\]
In particular, choosing $\eta =\Theta\!\bigl(1/\sqrt{T}\bigr)$ yields
\[
\min_{0\le t<T}\E\!\bigl[\|\nabla f(w_t)\|^2\bigr]=
\mathcal{O}\!\bigl(T^{-1/2}\bigr).
\]
\end{theorem}

\section*{Theoretical Properties}
\begin{itemize}[leftmargin=*]
\item \textbf{Stability.}  The bound \eqref{7} holds for any
non‑convex $f$ satisfying the coordinate‑wise smoothness; no convexity
is required.
\item \textbf{Probability weighting.}  The scaling $1/p_{i_t}$ in \eqref{1}
ensures $\E[g_t\mid w_t]=\nabla f(w_t)$.  Consequently the convergence rate
depends only on the *average* smoothness constant $L_{\max}:=\max_i L_i$
and not on the sampling distribution, apart from the variance term
$\sigma^2$ which may increase if extremely small $p_i$ are used.
\item \textbf{Comparison to uniform CD.}  Setting $p_i=1/d$ recovers the
standard randomized coordinate descent bound; with non‑uniform $p_i$
the bound \eqref{7} is unchanged but the variance $\sigma^2$ may be
smaller when more “important’’ coordinates are sampled more frequently.
\item \textbf{Hyperparameter sensitivity.}  The stepsize must respect
\eqref{6}.  If $\eta$ is taken larger, the term $L_i\eta$ may exceed $1$
and the descent lemma (Lemma~\ref{lem:descent}) fails, potentially leading
to divergence.
\end{itemize}

\section*{Discussion}
The theorem shows that, despite non‑convexity, the algorithm finds an
\emph{approximate stationary point} at the usual $O(T^{-1/2})$ rate.
The proof relies only on coordinate smoothness and the unbiasedness of
the scaled gradient \eqref{1}.  The presence of the factor $1/p_i$ is
crucial: without it the expectation of $g_t$ would be $p_i\nabla_i f$, and
the descent argument would involve an extra $p_i$ factor, degrading the
rate.  The bound is tight up to constants for stochastic first‑order
methods on smooth non‑convex functions.

\newpage

\section*{Preliminaries (Definitions and Lemmas)}
\begin{lemma}[Descent Lemma for a single coordinate]\label{lem:descent}
If $f$ satisfies Assumption~\ref{ass:smooth}, then for any $w\in\R^{d}$
and any scalar $\alpha$,
\[
f\bigl(w+\alpha e_i\bigr)
\le f(w)+\alpha\,\nabla_i f(w)+\frac{L_i}{2}\alpha^{2}.
\tag{8}
\]
\end{lemma}
\begin{proof}
Apply the fundamental theorem of calculus to the one‑dimensional
function $\phi(\tau)=f(w+\tau e_i)$ and use the Lipschitz bound
\eqref{4}.  The result is the standard quadratic upper bound.  ∎
\end{proof}

\section*{Main Proof (Step‑by‑Step Derivation)}
We prove Theorem~\ref{thm:main}.  Throughout the proof $w_t$ denotes the
random iterate generated by the algorithm, and expectations are taken
with respect to all randomness up to iteration $t$.

\bigskip
\noindent\textbf{Step 1.}  Apply Lemma~\ref{lem:descent} with
$\alpha=-\eta\,g_t^{(i)}/p_i$ for the sampled coordinate $i=i_t$.
Since all other coordinates stay unchanged,
\begin{align}
f(w_{t+1})
&\le f(w_t)-\eta\,\frac{1}{p_{i_t}}\nabla_{i_t} f(w_t)^{2}
      +\frac{L_{i_t}}{2}\Bigl(\frac{\eta}{p_{i_t}}\Bigr)^{2}
        \bigl(\nabla_{i_t} f(w_t)\bigr)^{2}.
\tag{9}
\end{align}
[Step 1]

\noindent\textbf{Step 2.}  Rearrange \eqref{9}:
\[
f(w_{t+1})\le f(w_t)-\frac{\eta}{2p_{i_t}}
\bigl(\nabla_{i_t} f(w_t)\bigr)^{2}
\;+\;
\frac{\eta L_{i_t}}{2p_{i_t}^{2}}
\bigl(\nabla_{i_t} f(w_t)\bigr)^{2}
\Bigl(\frac{\eta L_{i_t}}{p_{i_t}}-\!1\Bigr).
\tag{10}
\]
Because of the stepsize condition \eqref{6}, $\eta L_{i_t}\le1$,
hence the last term is non‑positive and can be dropped:
\[
f(w_{t+1})\le f(w_t)-\frac{\eta}{2p_{i_t}}
\bigl(\nabla_{i_t} f(w_t)\bigr)^{2}.
\tag{11}
\]
[Step 2]

\noindent\textbf{Step 3.}  Take conditional expectation w.r.t. $i_t$ given
$w_t$.  Using $\Pr(i_t=i)=p_i$,
\[
\E\!\bigl[f(w_{t+1})\mid w_t\bigr]
\le f(w_t)-\frac{\eta}{2}\sum_{i=1}^{d}
\bigl(\nabla_i f(w_t)\bigr)^{2}
= f(w_t)-\frac{\eta}{2}\,\|\nabla f(w_t)\|^{2}.
\tag{12}
\]
[Step 3]

\noindent\textbf{Step 4.}  Unroll the recursion \eqref{12} for $t=0,\dots ,T-1$
and take full expectation:
\[
\E\bigl[f(w_T)\bigr]
\le f(w_0)-\frac{\eta}{2}\sum_{t=0}^{T-1}
\E\!\bigl[\|\nabla f(w_t)\|^{2}\bigr].
\tag{13}
\]
[Step 4]

\noindent\textbf{Step 5.}  Use the lower bound $f(w_T)\ge f^{\inf}$ to obtain
\[
\frac{1}{T}\sum_{t=0}^{T-1}\E\!\bigl[\|\nabla f(w_t)\|^{2}\bigr]
\le\frac{2\bigl(f(w_0)-f^{\inf}\bigr)}{\eta T}.
\tag{14}
\]
[Step 5]

\noindent\textbf{Step 6.}  So far we have ignored the stochastic variance.
Recall that $g_t$ defined in \eqref{1} satisfies $\E[g_t\mid w_t]=\nabla f(w_t)$.
Write the update as $w_{t+1}=w_t-\eta\,\nabla f(w_t)-\eta\,(g_t-\nabla f(w_t))$.
Applying Lemma~\ref{lem:descent} with the full (unbiased) gradient
$\nabla f(w_t)$ and adding the error term yields
\[
f(w_{t+1})\le f(w_t)-\eta\|\nabla f(w_t)\|^{2}
+\frac{\eta^{2}L_{\max}}{2}\|g_t-\nabla f(w_t)\|^{2},
\tag{15}
\]
where $L_{\max}:=\max_i L_i$.
[Step 6]

\noindent\textbf{Step 7.}  Take expectation conditioned on $w_t$ and invoke
Assumption~\ref{ass:var}:
\[
\E\!\bigl[f(w_{t+1})\mid w_t\bigr]
\le f(w_t)-\eta\|\nabla f(w_t)\|^{2}
+\frac{\eta^{2}L_{\max}}{2}\sigma^{2}.
\tag{16}
\]
[Step 7]

\noindent\textbf{Step 8.}  Unrolling \eqref{16} and averaging over $T$ iterations
gives
\[
\frac{1}{T}\sum_{t=0}^{T-1}\E\!\bigl[\|\nabla f(w_t)\|^{2}\bigr]
\le\frac{f(w_0)-f^{\inf}}{\eta T}
+\frac{\eta L_{\max}}{2}\sigma^{2}.
\tag{17}
\]
[Step 8]

\noindent\textbf{Step 9.}  Since $L_{\max}\le 2$ under the stepsize choice
$\eta\le1/L_{\max}$, we can replace $L_{\max}$ by $2$ in \eqref{17},
which yields the cleaner bound
\[
\frac{1}{T}\sum_{t=0}^{T-1}\E\!\bigl[\|\nabla f(w_t)\|^{2}\bigr]
\le\frac{2\bigl(f(w_0)-f^{\inf}\bigr)}{\eta T}+2\eta\sigma^{2}.
\tag{18}
\]
[Step 9]

\noindent\textbf{Step 10.}  Inequality \eqref{18} is exactly the statement
\eqref{7} of Theorem~\ref{thm:main}.  ∎

\section*{Rate Derivation (With Constants)}
From \eqref{18} set $\eta = \sqrt{\frac{2\bigl(f(w_0)-f^{\inf}\bigr)}{\sigma^{2}T}}$
(which respects \eqref{6} for sufficiently large $T$).  Substituting,
\[
\frac{1}{T}\sum_{t=0}^{T-1}\E\!\bigl[\|\nabla f(w_t)\|^{2}\bigr]
\le 2\sqrt{\frac{2\sigma^{2}\bigl(f(w_0)-f^{\inf}\bigr)}{T}}
= \mathcal{O}\!\bigl(T^{-1/2}\bigr).
\tag{19}
\]
Thus the algorithm attains the optimal $O(T^{-1/2})$ rate for stochastic
first‑order methods on smooth non‑convex functions.

\section*{Special Cases}
\begin{enumerate}[leftmargin=*]
\item \textbf{Uniform sampling.}  If $p_i=1/d$ for all $i$, the scaled
gradient becomes $g_t=d\,\nabla_{i_t}f(w_t)e_{i_t}$.  The variance term
$\sigma^{2}$ typically decreases because each coordinate is sampled
equally often.

\item \textbf{Importance sampling.}  Choosing $p_i\propto L_i$ (or
$\propto |\nabla_i f(w_t)|$) reduces the expected squared scaling
$1/p_i^{2}$ and can dramatically lower $\sigma^{2}$, improving the
constant in \eqref{18} while leaving the $O(T^{-1/2})$ rate unchanged.

\item \textbf{Deterministic cyclic CD.}  When $p_i=1$ for a single
coordinate that cycles deterministically, the analysis above no longer
holds because the unbiasedness property \eqref{1} fails; the algorithm
may still converge but requires a different proof technique.
\end{enumerate}

\end{document}